\section{Introduction}
\label{sec:introduction}

Submarine cables form the physical backbone of the global Internet, carrying
nearly all intercontinental traffic and connecting users to services worldwide
~\cite{telegeography_submarine_map}. Cable disruptions can isolate regions,
degrade remote-service access, and redirect paths toward longer or congested
alternatives. Internet measurements describe the paths through logical
indicators such as routers, Autonomous Systems (ASes), and service replicas.

Logical and physical diversity describe different structures. Paths traversing
different routers or ASes may leave a country or region through the same
coastal area or converge on the same submarine corridor. Hereafter, country
refers to countries or regions. A service can therefore present a broad
network-layer path population while its feasible physical support occupies
only a few coastal directions. Joint analysis of the application target,
observed network path, and physical candidate space reveals whether
network-layer diversity is preserved across coastal directions.

Prior work measures user-path involvement with submarine infrastructure
~\cite{liu2020outofsight}, associates AS links with cable systems
~\cite{ye2023impact,ramanathan2023nautilus,wang2025calypso}, and constructs
cross-layer infrastructure maps~\cite{durairajan2013atlas,anderson2022igdb}.
Our paired distribution audit compares the network transitions and feasible
landing-region corridors of each country-service population. Both
distributions are constructed from the same atomic segments.

A \emph{landing-region corridor} is a direction-independent pair of bounded
landing regions connected by at least one active cable candidate. Router
locations, cable landing points and lifecycles, and propagation time determine
each segment's feasible corridor set. We allocate equal shares of segment
observation mass across this set and aggregate them into a candidate-support
distribution aligned with the network-transition distribution.

We organize the audit around three research questions:

\textbf{RQ1} How frequently do country-service paths contain a feasible inter-region submarine corridor?

\textbf{RQ2} How is observation mass distributed across feasible corridors?

\textbf{RQ3} Which country-service populations have broad network
transition distributions that project onto narrower corridor distributions?

We analyze 490,911 valid traceroutes from 18 public RIPE Atlas measurements
covering all 13 DNS Roots, Wikipedia, Reddit, Netflix Assets, and two
multi-target topology references. The aligned one-hour snapshot yields
3,291,063 atomic segments. Of these, 2,432,559 are projection-eligible and
124,350 support at least one feasible inter-region corridor. Under the 30-km
landing-region abstraction, 42.1\% of candidate-bearing segments support
multiple corridors.

We compare the distributions within each auditable country-measurement unit.
Ordered endpoint-AS pairs form the network distribution, and uniformly
allocated segment mass forms the corridor distribution. Top-2 share measures
absolute concentration, effective category count measures breadth, and
normalized entropy measures support-relative evenness. A common 80\% Top-2
threshold supplies a descriptive four-class view.

The 370-unit analysis identifies both contraction and broad support across
layers. It contains 222 DNS Root, 53 application, and 95 topology-reference
units. Median corridor Top-2 shares are 72.0\%, 70.3\%, and 30.6\%, with
effective corridor counts of 3.84, 3.51, and 18.01. Among 275 service-facing
units, 81 (29.5\%) combine broad network and concentrated corridor
distributions, and 158 (57.5\%) remain broad at both layers; 4 of 95
topology-reference units enter the first class.

Country-matched comparisons further show lower inter-region candidate
exposure for service-facing measurements. Island and archipelagic countries
form the highest-exposure group.

This paper makes three contributions. First, it formulates application-aware
submarine-cable measurement as a three-layer audit connecting services,
network transitions, and feasible physical corridors. Second, it develops a
reproducible projection that applies geographic, lifecycle, and propagation
feasibility and allocates observation mass uniformly over each segment's
corridor set. Third, it measures exposure, concentration, and cross-layer
contraction across DNS Roots, applications, and topology references,
including geographic and resolution sensitivity analyses.